\documentclass[11pt,a4paper]{article}
\usepackage[T1]{fontenc}
\usepackage{lmodern,amsmath,amssymb,amsthm,mathtools,booktabs,microtype}
\usepackage[margin=27mm]{geometry}
\usepackage[colorlinks=true,linkcolor=blue,urlcolor=blue]{hyperref}
\setlength{\emergencystretch}{2em}
\newtheorem{theorem}{Theorem}[section]
\newtheorem{lemma}[theorem]{Lemma}
\newtheorem{corollary}[theorem]{Corollary}
\newcommand{\Q}{\mathbb Q}\newcommand{\Z}{\mathbb Z}
\DeclareMathOperator{\rad}{rad}
\title{Rational height bounds for a signed power and a square}
\author{}\date{17 November 2026}
\begin{document}\maketitle

For a rational number $r=A/D$ in lowest terms, with $D>0$, write
$h(r)=\log\max(|A|,D)$; in particular $h(0)=0$.
Fix $\kappa>1$ and $K_\kappa\ge1$, and assume
\begin{equation}\label{eq:abc}
 c\le K_\kappa\rad(abc)^\kappa
\end{equation}
for every positive coprime integer triple $a+b=c$.

\begin{theorem}\label{thm:height}
Let $d$ be a positive even integer with
\[
 \epsilon=d-\kappa(d/2+2)>0,
 \qquad
 Q_{\kappa,d}=\frac{2\kappa-1}{\epsilon}
              =\frac{4\kappa-2}{(2-\kappa)d-4\kappa}.
\]
For every $x,y\in\Q$ and $\sigma\in\{1,-1\}$ with
$u=x^d+\sigma y^2\ne0$, one has
\begin{equation}\label{eq:bound}
 h(x)\le Q_{\kappa,d}h(u)+\frac{\log K_\kappa}{\epsilon}.
\end{equation}
No integrality, denominator, or real-sign assumption is imposed on
$x,y$, or on the nonzero number $u$.
\end{theorem}

\begin{proof}
If $x=0$, the result is immediate. If $y=0$, then
$h(x)=h(u)/d$; the inequality follows from
$Q_{\kappa,d}\ge1/d$ and $K_\kappa\ge1$.
Henceforth both numbers are nonzero.

Put $m=d/2$ and write
\[
 x=A/D,\qquad y=B/E,\qquad u=C/F
\]
in lowest terms, with $D,E,F>0$, $A,B,C\ne0$.
Let
\[
 \ell=\operatorname{lcm}(D^m,E),\qquad
 M=\ell/D^m,\qquad N=\ell/E.
\]
Clearing the common square denominator gives the integer identity
\begin{equation}\label{eq:triple}
 A^dM^2+\sigma B^2N^2=C\ell^2/F.
\end{equation}
In particular $F\mid\ell^2$.

Let
\[
 g=\gcd(A^dM^2,B^2N^2)=G^2,
 \qquad G=\gcd(|A|^mM,|B|N).
\]
We first verify the denominator facts needed below.
For every prime $p\mid\ell$, at least one term on the left of
\eqref{eq:triple} is a $p$-adic unit. Indeed, if
$m v_p(D)\ge v_p(E)$, then $v_p(M)=0$ and $A$ is a unit;
otherwise $v_p(N)=0$ and $B$ is a unit. Thus
\begin{equation}\label{eq:coprime}
 \gcd(g,\ell)=1,\qquad g\mid C,\qquad G\mid B.
\end{equation}
The divisibility $g\mid C$ follows from \eqref{eq:triple}, since
$\ell^2/F$ is an integer with all its prime factors in $\ell$.
For $G\mid B$, use $G\mid BN$ and $\gcd(G,N)=1$.

There is also the useful divisibility
\begin{equation}\label{eq:M}
 M^2\mid F.
\end{equation}
To prove it, let $p\mid M$, write $r=v_p(D)$ and $e=v_p(E)$,
and note $e>mr$. The term $y^2$ then has strictly smaller
valuation than $x^d$, namely $-2e$. Hence $v_p(F)=2e$, whereas
$v_p(M)=e-mr\le e$. This proves \eqref{eq:M} at every prime.
Finally the prime support of $\ell=D^mM$ gives
\begin{equation}\label{eq:radell}
 \rad(\ell)\le D\rad(M).
\end{equation}

Let $T$ be the largest of the absolute values of the three nonzero
terms in \eqref{eq:triple}. After division by $g$ and a possible
reordering, they form a positive coprime $abc$ triple. Its radical
$\mathcal R$ satisfies
\begin{equation}\label{eq:radical}
 \mathcal R
 \le |A|\,\frac{|B|}{G}\,\frac{|C|}{g}\,\rad(\ell).
\end{equation}
Indeed, the prime factors of the first divided term lie in $AM$,
those of the second lie in $(B/G)N$, and those of the third lie
in $(C/g)\ell$. All expressions used here are integers by
\eqref{eq:coprime}. Combining their prime supports proves the bound.

Define positive real numbers
\[
 H=(T/M^2)^{1/d},\qquad
 R=\max\{1,(F/|C|)^{1/d}\}.
\]
The three terms defining $T$ give
\[
 |A|\le H,\qquad |B|\le H^mM/N,\qquad D\le HR.
\]
Consequently $\max(|A|,D)\le HR$, and
\eqref{eq:radical}--\eqref{eq:radell} imply
\[
 \mathcal R\le
 |C|G^{-3} H^{m+2}R\,\frac{M\rad(M)}N.
\]
Apply \eqref{eq:abc} to $T/g$ and substitute $T=H^dM^2$:
\[
 H^\epsilon\le
 K_\kappa |C|^\kappa R^\kappa
 G^{2-3\kappa}M^{\kappa-2}\rad(M)^\kappa N^{-\kappa}.
\]
Since $\kappa>1$, one has $G^{2-3\kappa}\le1$ and
$N^{-\kappa}\le1$. Since $\rad(M)\le M$, and using
\eqref{eq:M}, one obtains
\begin{equation}\label{eq:H}
 H^\epsilon\le
 K_\kappa |C|^\kappa F^{\kappa-1}R^\kappa.
\end{equation}

Set $c=\log|C|$, $f=\log F$, and
$r=\log R=(f-c)_+/d$. Both $c,f$ are nonnegative, and
$h(u)=\max(c,f)$. From $h(x)\le\log H+r$ and
\eqref{eq:H},
\begin{equation}\label{eq:finalnumerator}
 \epsilon h(x)\le
 \log K_\kappa+\kappa c+(\kappa-1)f+(\kappa+\epsilon)r.
\end{equation}
If $c\ge f$, the last three terms are at most
$(2\kappa-1)c$. If $f\ge c$, substitute $r=(f-c)/d$;
the coefficients of $c$ and $f$ become
\[
 \kappa-\frac{\kappa+\epsilon}{d}
   =\frac{3\kappa}{2}-1+\frac\kappa d>0,
 \qquad
 \kappa-1+\frac{\kappa+\epsilon}{d}
   =\frac\kappa2-\frac\kappa d>0.
\]
Their sum is $2\kappa-1$. Therefore the same terms are at most
$(2\kappa-1)f$. Dividing \eqref{eq:finalnumerator} by
$\epsilon$ proves \eqref{eq:bound}.
\end{proof}

\begin{corollary}\label{cor:powered}
Under the same hypotheses, put $a=x^{d/2}$. Then
\[
 h(a)\le
 \frac{d(2\kappa-1)}{(2-\kappa)d-4\kappa}\,h(u)
 +\frac{d\log K_\kappa}{2\epsilon}.
\]
For fixed $1<\kappa<2$, the displayed height coefficient tends to
$(2\kappa-1)/(2-\kappa)$ as $d$ tends to infinity through the
admissible even integers.
\end{corollary}
\begin{proof}
Use $h(x^{d/2})=(d/2)h(x)$ in Theorem~\ref{thm:height}.
\end{proof}

The weaker coefficient
\[
 C_{\kappa,d}=
 \frac{(4+\kappa)d-6\kappa}{(2-\kappa)d-4\kappa}
\]
from the earlier plus-square argument also holds for the signed
equation, since $C_{\kappa,d}>Q_{\kappa,d}$ throughout the stated
range. The square gcd and minimal common denominator in the proof
avoid the height cost of multiplying by the full $d$th power of
the denominator of $u$.

The exclusion $u=0$ is necessary: one may set
$y=x^{d/2}$ in the minus case with arbitrarily large $h(x)$.
The theorem concerns one square. It gives no assertion for a sum or
difference containing two independently varying squares, and no
prime-pattern realization or defining-variable count follows from
the height bound alone.

\begin{thebibliography}{9}\raggedright
\bibitem{Single}\emph{A single-cover criterion from a mixed-power height
inequality}, repository manuscript,
\nolinkurl{papers/pole_parameters/single_cover_criterion_20261117_v56.tex}.
\end{thebibliography}
\end{document}
